\section{The mod-four Bockstein and the exact lattice}\label{sec:integral-lattice}

The signed cut determines whether an exterior coordinate extends,
but an integral extension may have a fractional \(K_{\Q}\)-coordinate.
We determine that coordinate from the torsion boundary in \(H_2(J)\).
Write
\(\beta_{\rm Bock}:H_3(J;\Z/k)\to H_2(J)\) for the coefficient Bockstein,
using the same notation on restricted interfaces below.
Our convention is
\[
 \beta_{\rm Bock}[\overline C]=[\partial C/k]
\]
when \(C\) is an integral lift of a mod-\(k\) cycle
\(\overline C\). Under \eqref{eq:tube-11},
the signed root boundary views
a class of \(H_3(W;\Z/k)\) in \(D/kD\).
Let \(i=i_W:J\hookrightarrow W\) be the inclusion. We also write
\(\pi^!\) for the induced map \(H/kH\to D/kD\);
it is well defined because \(\pi^!(kH)\subseteq kD\).

\begin{lemma}[Bockstein equals finite-map transfer]\label{lem:bock}
For \(b\in H/kH=H_2(J)\), and any
\(\eta\in H_3(J;\Z/k)\) with \(\beta_{\rm Bock}\eta=b\),
\begin{equation}
\Psi(i_*\eta)=\pi^!b\pmod{kD},
\label{eq:tube-14}
\end{equation}
with the orientation of \(\Psi\) fixed after \eqref{eq:tube-12}.
\end{lemma}
\begin{proof}
First note that the ambiguity in \(\eta\) is reduction of
an integral class in \(H_3(J)=\Z\), and this class dies
in \(H_3(W)\) by Lemma~\ref{lem:mv}; thus the left-hand
side is unambiguous. It suffices to test on loops in
\(\Hh\) avoiding the finite branch discriminant.
Take \(w_0=\lambda/2\), with every root in
\(|w|<\lambda/4\). On the \(\theta=0\), \(h=c\)
model of the interface lens fiber, for \(p\in\Hh_c\)
the \(y\)-angle loop
\[
\gamma_p
=\{w=w_0,\ y=M\lambda^\alpha e^{i\phi},\
z=R_s(p,w_0)/y\}
\]
is the positive generator of \(H_1(L)=\Z/k\).
Indeed, at \(s=0\) the quotient
\((U,V)\mapsto(U^k,V^k,UV)\) lifts one positive turn of the
\(y\)-angle to a path ending at the deck transform
\((U,V)\mapsto(\omega U,\omega^{-1}V)\), where
\(\omega=e^{2\pi i/k}\). This generates the deck group.
The anchored \(s\)-continuation preserves the same oriented loop.
Let \(\widetilde\gamma_p\) denote its image on the actual
\(f_s=c\) interface under the explicit lift in
Lemma~\ref{lem:continuations}. That lift preserves the
\(h\)-page label \(p\), \(w_0\), and the exact \(y\)-angle,
though physical \(p\) moves. We compute below in the
\(\theta=0\) model and transport all relative classes to
the actual \(J\) and \(W\) by this isotopy. In the model,
the exact sequence of the pair \((L,\gamma_p)\) reads
\[
0=H_2(L)\to H_2(L,\gamma_p)
\xrightarrow{\partial}H_1(\gamma_p)=\Z
\to H_1(L)=\Z/k.
\]
It supplies a unique relative class \(S_p\) with
\(\partial S_p=k\gamma_p\).

Let \(F_p\) be the inner \(yz=R_s(p,w)\) fiber above a
regular \(p\), and let \(r_1,\dots,r_k\) be the roots in
the \(w\)-disk \(\Delta\). Let
\(X_p=X(\Delta,\{r_1,\dots,r_k\})\) be its balanced collapse core,
and let \(\gamma_p^{\rm core}\) be the circle over \(w_0\) in that
core. The balancing retraction does not fix the outer circle
\(\gamma_p\) pointwise. It gives a map of pairs
\[
 r_p:(F_p,\gamma_p)\longrightarrow(X_p,\gamma_p^{\rm core})
\]
that is an absolute homotopy equivalence and restricts to the
orientation-preserving circle homeomorphism retaining the \(y\)-angle.
The long exact sequences of the two pairs therefore show that
\(r_p\) induces an isomorphism on relative homology.
The constant-angle section of \(X_p\) is a disk meeting
\(\gamma_p^{\rm core}\) in one point. Adjoining that disk to the
relative subspace does not change relative homology; quotienting
it and \(\gamma_p^{\rm core}\) gives the circle suspension of
\(\Delta/(\{r_1,\dots,r_k,w_0\})\). Hence there is a natural isomorphism
\begin{equation}
\begin{split}
 H_2(F_p,\gamma_p)
 &\cong H_1(\Delta,\{r_1,\dots,r_k,w_0\})\\
 &\cong\{(a_1,\dots,a_k,a_0)\in\Z^{k+1}:
             a_1+\cdots+a_k+a_0=0\}.
\end{split}
\label{eq:tube-15}
\end{equation}
Use \(r_p\) to identify the oriented boundary classes of
\(\gamma_p\) and \(\gamma_p^{\rm core}\) in this calculation.
The class of an arc from \(w_0\) to \(r_i\), swept by
the \(y\)-angle, has root coefficient \(+1\) and
\(\gamma_p\)-boundary \(+1\), fixing the convention.
The image of \(S_p\) has \(a_0=-k\).
To compute this one fiber class, let \(r_i(t)\), \(0\le t\le1\),
be any braid of distinct roots in \(|w|<\lambda/4\), starting at
the roots of \(R_s(p,\cdot)\) and returning to the same unordered
root set. Put
\(Q_t(w)=\prod_{i=1}^k(w-r_i(t))\), so \(Q_1=Q_0\), and consider
the compact fibers \(F_t=\{yz=Q_t(w)\}\cap Y_\lambda\), their
boundaries \(L_t\), and the circles \(\gamma_t\) defined by the
same \(w_0\) and \(y\)-angle formula as \(\gamma_p\), with
\(R_s(p,w_0)\) replaced by \(Q_t(w_0)\).
Choose a smooth disk velocity \(v_t(w)\), supported in
\(|w|<\lambda/4\), equal to \(\dot r_i(t)\) near each moving root.
Then
\[
 \ell_t(w)=\frac{\dot Q_t(w)+Q'_t(w)v_t(w)}{Q_t(w)}
\]
extends smoothly across the roots. Indeed, near \(r_i(t)\), writing
\(Q_t=(w-r_i)B_{i,t}\) gives
\(\ell_t=(\dot B_{i,t}+\dot r_i B'_{i,t})/B_{i,t}\), with
\(B_{i,t}\ne0\).

The \(y\)- and \(z\)-faces of these fibers never meet, since
\(|Q_t(w)|\le(5\lambda/4)^4<M^2\lambda^4\).
Choose a smooth real cutoff \(\chi\), depending on
\(|y|^2,|z|^2\), that is zero near the \(y\)-face and one near the
\(z\)-face. The time-dependent lift
\[
 \dot w=v_t(w),\qquad
 \dot y=\chi\operatorname{Re}(\ell_t)y,\qquad
 \dot z=\bigl(\ell_t-\chi\operatorname{Re}(\ell_t)\bigr)z
\]
is smooth also at the roots and axes, and differentiation gives
\(\frac{d}{dt}(yz-Q_t(w))=\ell_t(yz-Q_t(w))\).
It is tangent to the \(y\)-face because \(\dot y=0\) there, to the
\(z\)-face because the coefficient of \(z\) is purely imaginary
there, and to the \(w\)-face because \(v_t=0\) there; these statements
also cover all corners. Compactness therefore gives transport of
the full triples \((F_t,L_t,\gamma_t)\). On \(\gamma_t\), it fixes
\(w_0\) and \(y=M\lambda^\alpha e^{i\phi}\) exactly, and gives
\(\dot z=\dot Q_t(w_0)/y\). Thus the return fixes the parametrized
circle \(\gamma_p\) pointwise.

The return maps of \((L,\gamma_p)\) and \((F_p,\gamma_p)\) commute
with inclusion. The former fixes the unique class \(S_p\) with
boundary \(k\gamma_p\); hence the latter fixes its image in
\(H_2(F_p,\gamma_p)\). To compute that same return map in
\eqref{eq:tube-15}, observe that the lift covers the disk flow
\(\psi_t\) of \(v_t\) and preserves the \(y\)-angle wherever defined.
It therefore intertwines the balancing collapses with the map
\((w,e^{i\phi})\mapsto(\psi_t(w),e^{i\phi})\) on the circle models;
collapsed root fibers are carried to collapsed root fibers.
The endpoint isomorphism in
\eqref{eq:tube-15} consequently sends its action to the permutation
of the root coefficients \(a_i\), with \(a_0\) fixed. Every root
permutation is available, and no additional term can occur because
the endpoint map is an isomorphism. Invariance of the image of
\(S_p\), together with \(a_0=-k\), now gives
\(a_1=\cdots=a_k=1\). Reducing \(S_p\) modulo \(k\)
closes its boundary and yields the diagonal root vector in the
absolute mod-\(k\) collapse homology.

Now work over an oriented base loop \(a\subset\Hh_c\) missing
branch values, still in the \(\theta=0\), \(h=c\) model. Write
\(\D_a=\pi^{-1}(a)\). Let \(W_a\) be the restriction of the inner
family \(F_p\) to \(p\in a\), and let \(J_a\subset W_a\) be its
restricted lens-interface family, with inclusion
\(i_a:J_a\hookrightarrow W_a\). The compatible isotopy of
Lemma~\ref{lem:continuations} transports their inclusions to the actual
pieces \(J\subset W\); the vertical maps below use this transport.
The same circle-first suspension and negative relative root boundary
define
\[
\Psi_a:H_3(W_a;\Z/k)\longrightarrow H_1(\D_a;\Z/k).
\]
The target remains the root homology in the \(h\)-model, with its
inclusion into \(H_1(\D;\Z/k)\). The actual covering monodromy
over \(a\) may be a smaller braid subgroup, with several orbits.
Over a small arc
\(I\subset a\), the same transport tracks roots and preserves the
circle angle. Collapse relative to the constant-angle
section, take the relative root boundary, and slice
over \((I,\partial I)\). Each operation is natural
under this product trivialization: on chains its
relative map is the fiber operation times
\([I,\partial I]\). Restriction and inclusion give
the following naturality diagram:
\begin{equation}
\begin{array}{ccccc}
H_3(J_a;\Z/k)&\xrightarrow{i_{a*}}&H_3(W_a;\Z/k)
 &\xrightarrow{\Psi_a}&H_1(\D_a;\Z/k)\\
\downarrow&&\downarrow&&\downarrow\\
H_3(J;\Z/k)&\xrightarrow{i_*}&H_3(W;\Z/k)
 &\xrightarrow{\Psi}&H_1(\D;\Z/k).
\end{array}
\label{eq:tube-16}
\end{equation}
The common orientation of the normal circle fixes the
sign of the two right-hand maps. The fiber endpoint
calculation gives the diagonal root vector after
slicing the top row at \(p\).
We construct a Bockstein lift on \(J_a\) by first closing a swept
chain at the seam. The anchored circles trace the base-first torus
\(a\times\gamma_p\) over \(a\). Cut \(a\) at one point, choose a
representative
\(S_0\) of \(S_p\) with \(\partial S_0=k\gamma_p\), transport it
along the resulting interval, and call its endpoint \(S_1\).
The anchored circle is fixed at the seam, so
\(\partial S_0=\partial S_1=k\gamma_p\). Uniqueness of the relative
class makes \(S_1-S_0\) a null-homologous \(2\)-cycle in \(L\);
choose a seam cap \(C_a\) with
\(\partial C_a=S_1-S_0\), using \(H_2(L)=0\).
Let \(T_a\) be the interval-first sweep. Then
\[
 \partial T_a=S_1-S_0-k(a\times\gamma_p),\qquad
 \partial(T_a-C_a)=-k(a\times\gamma_p).
\]
The second boundary is divisible by \(k\), so the reduction is a cycle.
Set
\[
\eta_a=-[(T_a-C_a)\bmod k]\in H_3(J_a;\Z/k),
\]
and let \(\eta\) be its image in \(H_3(J;\Z/k)\) under inclusion
and transport.
The Bockstein convention gives
\(\beta_{\rm Bock}\eta_a=[a\times\gamma_p]\) on \(J_a\), and
\(\beta_{\rm Bock}\eta=[a\times\gamma_p]\) on \(J\), with the
torus transported in the second identity. Changing the cap changes
\(\eta_a\) by the reduction modulo \(k\) of an integral class from
\(H_3(L)\), and changes \(\eta\) by its image. The integral class
maps to zero in \(H_3(W)\) by the fiber
calculation used in Lemma~\ref{lem:mv}, so its mod-\(k\) image
dies as well.
Thus both
\(\Psi i_*\eta\) and the full transfer \(\pi^![a]\) are
images of classes in \(H_1(\D_a;\Z/k)\).
Choose the slice away from the seam. Slicing at \(p\in a\) sends each oriented component circle
of the unramified cover \(\D_a\to a\) to the sum of the
roots in its monodromy orbit. Distinct orbits have
disjoint supports, so this slice map is injective on
\(H_1(\D_a;\Z/k)\). The fiber calculation gives the
diagonal root vector for \(\Psi_a i_{a*}\eta_a\), and the
full transfer has the same vector. Their classes agree
already in \(H_1(\D_a;\Z/k)\), hence in \(D/kD\).
Such base loops generate
\(H_1(\Hh)\), proving \eqref{eq:tube-14}.
\end{proof}

We now define the rational coordinates of an integral page class.
For \(x\in H_3(\T_s)\), let \(z=\partial x\) and
\(b=(1\otimes\tau)z\in H/kH\). Since
\(kz\in H\otimes\ker\tau\), Lemma~\ref{lem:mv} supplies a unique
\(x_{\rm ext}(kz)\) in the exterior image with coordinate \(kz\).
Define
\[
u:=kx-x_{\rm ext}(kz)\in K,\qquad
x\longmapsto (u/k,z)\in K_\Q\oplus(H\otimes E)_\Q.
\]
We call \(u\) the integral numerator of the inner coordinate.
The Bockstein lemma supplies \(+\pi^!b\); the cap calculation below
shows that subtracting the interface cap gives the numerator residue
\(u\equiv-\pi^!b\pmod{kD}\).

\begin{lemma}[integral extension]\label{lem:lattice}
In these rational coordinates, the image of \(H_3(\T_s)\) is exactly
the lattice
\[
\Lambda_\pi=
\left\{(\pr l,z): l\in D,\ z\in H\otimes E,\quad
\pi_*l\bmod kH=(1\otimes\tau)z\right\}
\subset K_\Q\oplus(H\otimes E)_\Q
\]
from \eqref{eq:tube-1}, where
\(\pr l=l-k^{-1}\pi^!\pi_*l\).
\end{lemma}
\begin{proof}
Fix \(x\), with \(z,b,u\) as above. Orient \(J\) as the outgoing
boundary of \(W\), and represent \(x\)
by page chains \(w\in C_3(W)\) and
\(v\in C_3(M_{\rm ext})\) with
\(\partial w=b_J\) and \(\partial v=-b_J\).
The oriented prism calculation \eqref{eq:mv-cut-chain}--\eqref{eq:mv-strip}
gives \([b_J]=b=(1\otimes\tau)z\) in
\(H_2(J)=H/kH\). Because this
group is killed by \(k\), choose an
integral \(A\in C_3(J)\) with
\(\partial A=k b_J\). Write
\(i_{\rm ext}:J\hookrightarrow M_{\rm ext}\) for the exterior inclusion.
The chains
\begin{equation}
 w_k:=kw-i_{W*}A,\qquad
 v_k:=kv+i_{{\rm ext}*}A
\label{eq:tube-16a}
\end{equation}
are closed in their respective pieces.
Their sum in the full page is \(kx\):
the two copies of \(A\) cancel on \(J\).
The exterior class \([v_k]\) has
second coordinate \(kz\), hence is
\(x_{\rm ext}(kz)\) in the full page;
its possible \(H_3(J)\) ambiguity dies
after attaching \(W\). Consequently
\(u=[w_k]\in K\). Modulo \(k\),
\([w_k]=-i_{W*}[A\bmod k]\).
The coefficient convention gives
\(\beta_{\rm Bock}[A\bmod k]=[\partial A/k]=[b_J]=b\).
With the \(\Psi\) orientation fixed after
\eqref{eq:tube-12}, Lemma~\ref{lem:bock}
gives \(\Psi i_*[A\bmod k]=+\pi^!b\).
The preceding minus sign in
\([w_k]=-i_{W*}[A\bmod k]\) therefore gives
\begin{equation}
 u\equiv-\pi^!b\pmod{kD}.
\label{eq:tube-17}
\end{equation}
Choose \(b_H\in H\) lifting \(b\) and put
\(l=(u+\pi^!b_H)/k\). The congruence makes
\(l\in D\), and
\(\pi_*l=b_H\), \(\pr l=u/k\).
Conversely, if \(l,z\) obey the condition in
\eqref{eq:tube-1}, the last arrow of \eqref{eq:tube-13} kills \(z\).
Choose any \(x\) with that \(z\), and put
\[
u_x:=kx-x_{\rm ext}(kz)\in K,\qquad
u_l:=kl-\pi^!\pi_*l\in K.
\]
Here \(u_x\) is the integral numerator of the rational inner
coordinate \(u_x/k\).
By \eqref{eq:tube-17}, \(u_x-u_l\in K\cap kD\).
This intersection equals \(kK\):
if \(kv\in K\) for \(v\in D\), then
\(k\pi_*v=0\), and \(H\) is torsion free,
so \(v\in K\). Adjust \(x\) by
\((u_l-u_x)/k\in K\) through the first
arrow of \eqref{eq:tube-13}. The adjusted class has
rational coordinate \((\pr l,z)\).
Thus every pair in \eqref{eq:tube-1} occurs.
Replacing \(b_H\) by
\(b_H+kh\) changes \(l\) by \(\pi^!h\) and
does not change \(\pr l\), so the coordinate
description is independent of the lift.
\end{proof}

We have now located every integral page class in
\(K_{\Q}\oplus(H\otimes E)_{\Q}\). The residue condition in
\eqref{eq:tube-1} retains the particular projection \(\pi_*\);
an abstract isometry of the two plane Seifert lattices does not
identify it. It remains to compute the ordered filling form on
this exact lattice.
